\documentclass[10pt,a4paper]{amsart}
\usepackage[margin=19mm]{geometry}
\usepackage{amsmath,amssymb,mathtools}
\usepackage[scr=rsfs,cal=euler]{mathalfa}
\usepackage{xfrac}
\usepackage[unicode,hidelinks,bookmarks=false]{hyperref}
\newcommand{\ie}{i.e.\ }
\newcommand{\cf}{cf.\ }
\newcommand{\resp}{respectively\ }
\newcommand{\Subsection}{\subsection}
\providecommand{\nobreakdash}{\nobreak\hbox{-}}
\setlength{\emergencystretch}{2em}
\multlinegap0pt
\usepackage{bm}
\usepackage{bbm}
\usepackage{dsfont}
\usepackage{xspace}
\usepackage{cancel}
\usepackage{mathtools}
\usepackage[shortlabels]{enumitem}
\usepackage{amsthm}
\makeatletter
\@ifundefined{thm}{\newtheorem{thm}{Thm}}{}
\@ifundefined{lemma}{\newtheorem{lemma}[thm]{Lemma}}{}
\makeatother
\newcommand{\Col}[1]{\mathcal{C}(#1)}
\newcommand{\F}{{\mathbf F}}
\newcommand{\va}{\vec{a}}
\newcommand{\vn}{\vec{n}}
\newcommand{\vu}{\vec{u}}
\title[Clonoids over vector spaces]{Clonoids over vector spaces}
\author{Stefano Fioravanti, Michael Kompatscher, Bernardo Rossi}
\date{Source: arXiv:2602.04034v1 (CC BY 4.0)}
\begin{document}
\maketitle

\noindent\emph{Source and licence.} Clonoids over vector spaces. Version arXiv:2602.04034v1 (CC BY 4.0), distributed under the Creative Commons Attribution 4.0 International licence, as recorded in the arXiv licence metadata for this exact version and shown on the abstract page https://arxiv.org/abs/2602.04034v1. Full rights notice: licenses/arxiv-2602.04034-CC-BY-4.0.txt.\par\smallskip

\noindent\textit{Editorial scope.} Counted unit: the Lemma whose source label is \texttt{lemma:representation} in the article (source0.tex lines 1111--1120, source label \texttt{lemma:representation}), delivered together with the article's own complete original proof of it (source0.tex lines 1122--1165, one \texttt{proof} environment of 44 lines). No mathematics is written, repaired or re-derived: statement and proof are the article's own text line for line. Required context retained uncounted: none. Every definition, display and label of those lines is retained unchanged. Omitted from this package and not redistributed: 1--1110, 1121--1121, 1166--1542. Nothing inside a retained excerpt is omitted or altered: every retained body is delivered exactly as the article's lines are written, so no omission receipt of any kind accompanies this package. Citations: the retained excerpts carry no citation command at all, so no bibliography entry of the article is transcribed and no reference is redistributed. Reference adaptation: none was needed and none was made; every reference inside the delivered unit resolves inside it, so no unresolved reference or citation is produced. Printed slips kept literal: no printed word, symbol, hypothesis or number of the counted statement or of its proof was altered. Layout and class: the article's own class is \texttt{amsart} with packages including a4wide, amsmath, amsfonts, amssymb, amscd, graphicx, amsthm, cite, mathrsfs, url, hyperref, color, comment, enumitem; the wrapper is the lane's pinned \texttt{amsart} editorial wrapper at 10pt on A4, which restates the theorem-like environments the retained text needs and adds the article's own macro definitions that the retained text needs (\texttt{\textbackslash Col}, \texttt{\textbackslash F}, \texttt{\textbackslash va}, \texttt{\textbackslash vn}, \texttt{\textbackslash vu}), with extra wrapper packages (bm, bbm, dsfont, xspace, cancel, mathtools, enumitem). The delivered numbering is the unit's own. Assets: no retained span contains a figure environment, a graphics-inclusion command, a PDF-page inclusion, a table or a TikZ picture; no image file is copied into this package, the package declares no asset dependency and no image file is redistributed with it. Licence: version arXiv:2602.04034v1 of this article is distributed under the Creative Commons Attribution 4.0 International licence, as recorded in the arXiv licence metadata for this exact version and shown on the abstract page https://arxiv.org/abs/2602.04034v1. A full rights notice is delivered at licenses/arxiv-2602.04034-CC-BY-4.0.txt. Characters: every retained excerpt is pure ASCII, so the delivered engine, XeLaTeX with its default Latin Modern text font, reports no missing character.
\begin{lemma} \label{lemma:representation}
Let $X \in \mathfrak T_n$, and 
let $X = X_0 + AU = X_0 + \tilde A \tilde U$ 
be two admissible representations of it. Furthermore, let $X_i = X_0 + \sum_{j=1}^i \va_j\vu_j^T$ and $\tilde X_i = X_0 + \sum_{j=1}^i \tilde \va_j \tilde \vu_j^T$ for $i \in[n]_0$. Then:
\begin{enumerate}
\item\label{ite:lemma:representation2} $\Col{X_0} \cap \Col{X_1} \cap \ldots \cap \Col{X_i} = \Col{\tilde X_0} \cap \Col{\tilde X_1} \cap \ldots \cap \Col{\tilde X_i}$ for 
every $i\in[n]_0$.
\item\label{ite:lemma:representation1} there is an lower triangular matrix $T \in \F^{n\times n}$ with $A = \tilde A T^{-1}$ and $U = T \tilde U$. 
\end{enumerate}
\end{lemma}

\begin{proof}
Since $AU = \tilde A \tilde U$ are two rank factorizations of $X-X_0$ we already know that there exists an invertible matrix $T= (t_{ij})_{i,j =1}^n$ with $A = \tilde AT$ and $U = T^{-1} \tilde U$. 


% We are going to prove by induction on $i=0,1,\ldots,n$
% that
% \begin{enumerate}[label=(\alph*)]
% \item\label{ite:lemma:representationa} $\Col{X_0} \cap \Col{X_1} \cap \ldots \cap \Col{X_i} = \Col{\tilde X_0} \cap \Col{\tilde X_1} \cap \ldots \cap \Col{\tilde X_i}$, and that
% \item\label{ite:lemma:representationb} the first $i+1$ columns of $T$ look like a lower triangular matrix, i.e., $t_{ab} = 0$ for all $a < b$ with $b\leq i+1$.
% \end{enumerate}

% Note that \ref{ite:lemma:representationb} implies that the same statement holds also for the inverse of $T^{-1}$. This can be seen from the following formula, if we consider $T$ as a block matrix with blocks of dimension $i+1$ and $n-(i+1)$:

% \begin{equation}
% T = \begin{bmatrix} T_{11} & 0 \\ T_{21} & T_{22} \end{bmatrix} \Rightarrow T^{-1} = \begin{bmatrix} T_{11}^{-1} & 0 \\ -T_{22}^{-1}T_{21}T_{11}^{-1} & T_{22}^{-1} \end{bmatrix}.
% \end{equation}
For $i\in \{0,1,\ldots,n\}$, let $\vn_i, \tilde \vn_i$ denote vectors such that $\langle \vn_i \rangle = \Col{X_i}^\bot = \ker(X_i^T)$ and $\tilde \vn_i = \Col{\tilde X_i}^\bot = \ker(\tilde X_i^T)$. We are going to prove by induction on $i=0,1,2,\ldots, n$ that
\begin{enumerate}[label=(\alph*)]
\item\label{ite:lemma:representationa} $\langle \vn_0,\ldots,\vn_i \rangle = \langle \tilde \vn_0,\ldots,\tilde \vn_i \rangle$.
\item\label{ite:lemma:representationb} $\langle \vu_1,\ldots,\vu_i, \vu_{i+1} \rangle = \langle \tilde \vu_1,\ldots,\tilde \vu_{i}, \tilde \vu_{i+1} \rangle$.
\end{enumerate}

By considering the orthogonal spaces, we can see that \ref{ite:lemma:representationa} is equivalent to $\Col{X_0} \cap \Col{X_1} \cap \ldots \cap \Col{X_i} = \Col{\tilde X_0} \cap \Col{\tilde X_1} \cap \ldots \cap \Col{\tilde X_i}$. Furthermore, if $\langle \vu_1,\ldots, \vu_{l} \rangle = \langle \tilde \vu_1,\ldots,\tilde \vu_{l} \rangle$ holds for every $l\leq i+1$, then $TU = \tilde U$ implies that $t_{lj} = 0$ for all $l,j$ satisfying $i+1 \geq l>j$; in other words, $T$ looks like a lower triangular matrix on its first $i+1$ rows. So, if we prove (a) and (b) for $i=0,1,\ldots,n$, we are done.

We first prove \ref{ite:lemma:representationa} and \ref{ite:lemma:representationb} for $i=0$. By $X_0 = \tilde X_0$ clearly $\langle \vn_0 \rangle = \langle \tilde \vn_0 \rangle$ holds. For \ref{ite:lemma:representationb}, observe that $\vn_0^T\va_i \neq 0 \Leftrightarrow \va_i \not\in \Col{X_0} \Leftrightarrow i = 1$. This implies that $\vn_0^TX = \vn_0(X_0 + \sum_{j=1}^n \va_j\vu_j^T) = \vn_0^T\va_1\vu_1^T$. Symmetrically, we obtain $\vn_0^TX =\vn_0^T\tilde \va_1\tilde \vu_1^T$. Thus $\tilde \vu_1 = (\vn_0^T\va_1)^{-1}(\vn_0^T\va_1)\vu_1$, which finishes the proof of the base step.

Next, let us prove the induction step $1,\ldots,i-1\to i$. We start with \ref{ite:lemma:representationa}. By induction assumption $\langle \vn_0,\ldots,\vn_{i-1} \rangle = \langle \tilde \vn_0,\ldots,\tilde \vn_{i-1} \rangle$, thus, it is enough to show that $\tilde \vn_{i} \in \langle \vn_0,\ldots,\vn_{i-1}, \vn_i \rangle$. This is equivalent to finding coefficients $r_0,\ldots,r_{i-1}$, such that the equation $(\vn_i + \sum_{j=0}^{i-1}r_j\vn_j)^T \tilde X_i = 0$ holds.

By the induction assumption we know that $T$ has the shape of lower triangular matrix on its first $i$ rows. Recall that $\tilde A = AT$, and so $\tilde \va_j = \sum_{l\geq j} t_{lj} \va_l$ for every $j\leq i$. For $j < i$ this implies $\vn_j^T \tilde \va_{j+1} = t_{j+1,j+1} \vn_j^T \va_{j+1} \neq 0$. For all $j+1<l\leq i$ this further implies $\vn_j^T \tilde \va_{l} = 0$ holds. By both of these observations, for every $j < i$, we obtain 
\begin{equation} \label{eq:uis1} \vn_j^T \tilde X_i = \vn_j^T (\tilde X_j + \sum_{k>j} \tilde \va_{k}\tilde \vu_k^T) = c_j \vu_{j+1}^T, \text{ where } c_j = t_{j+1,j+1} \vn_j^T \va_{j+1} \neq 0.
\end{equation} On the other hand, note that
\begin{equation} \label{eq:uis2}
\vn_i^T \tilde X_i = \vn_i^T (\tilde X_i - X_i) = \vn_i^T (\sum_{j\leq i} \tilde \va_j \tilde\vu_j^T - \va_j\vu_j^T) \in \langle \vu_1^T, \ldots, \vu_j^T\rangle.
\end{equation}
The latter follows from $\langle \tilde\vu_1, \ldots, \tilde\vu_j\rangle = \langle \vu_1, \ldots, \vu_j\rangle$, which holds by the induction hypothesis. Thus, there are coefficients $d_1,\ldots, d_n$, such that $\vn_i^T \tilde X_i = \sum_{j\leq i} d_j \vu_j$. We conclude that $(\vn_i + \sum_{j=0}^{i-1}r_j \vn_j)^T \tilde X_i = 0$ is equivalent to the equation 
\begin{equation} \label{eq:uis3}
\sum_{j = 1}^i d_j \vu_j^T + \sum_{j=0}^{i-1} r_jc_j \vu_{j+1}^T = \sum_{j=0}^{i-1} (d_{j-1}+r_jc_j) \vu_{j+1}^T= 0 
\end{equation}
Equation \eqref{eq:uis3} has the solution $r_j = -d_{j-1}c_j^{-1}$ for $j=0,\ldots,i-1$ (recall that $c_j$ is invertible by \eqref{eq:uis1}). Thus $\tilde \vn_{i} \in \langle \vn_0,\ldots,\vn_{i-1}, \vn_i \rangle$, and we are done.

In order to prove \ref{ite:lemma:representationb} we are again going to use that we can write $\tilde \vn_i$ as linear combination $\tilde \vn_i = \sum_{j = 0}^{i} r_j \vn_j$ (with $r_i = 1$). By multiplying with $X$ we obtain $\tilde \vn_i^T X = \sum_{j=0}^i r_j \vn_j^T X = \sum_{j=0}^i r_j (\vn_j^T \va_{j+1})\vu_{j+1}^T$. On the other hand, clearly also $\tilde \vn_i^T X  = (\vn_i^T \tilde \va_{i+1}) \tilde \vu_{i+1}^T$ holds. Since $\vn_i^T \tilde \va_{i+1} \neq 0$, this implies that $\tilde \vu_{j+1} \in \langle \vu_{1}, \ldots \vu_{i+1}\rangle$. Since $\tilde \vu_{j+1} \notin \langle \tilde \vu_{1}, \ldots \tilde \vu_{i}\rangle = \langle \vu_{1}, \ldots \vu_{i}\rangle$, this implies that $\langle \tilde \vu_{1}, \ldots \tilde \vu_{i+1}\rangle = \langle \vu_{1}, \ldots \vu_{i+1}\rangle$, so we are done.

Since items \ref{ite:lemma:representationa} \ref{ite:lemma:representationb} hold for all $i=0,1,\ldots, n$, this finishes the proof of the lemma.
\end{proof}

\end{document}
